\section{Syntactic Analysis}\label{sec:syntax}

\subsection{What the grammar has to describe}

Real Yaound\'e speech does not arrive as tidy subject--verb--object sentences. The grammar
has to admit at least the following, all of which occur in the data:

\begin{itemize}[leftmargin=1.4em, itemsep=0.2em]
  \item clauses juxtaposed without any conjunction, separated only by a comma or by
        nothing at all (\fw{on est l\`a, \c ca va});
  \item a fronted prepositional phrase before the subject (\fw{dans le taxi, ...});
  \item bare interjections and frozen formulas standing alone as complete utterances;
  \item serial verbs -- two or more verbs in sequence with no linking word;
  \item pre-verbal negation and adverbs interleaved inside the verb group;
  \item subject dropping entirely, leaving a bare predicate.
\end{itemize}

\subsection{Original context-free grammar}

The grammar is authored in its natural, human-readable form -- including left recursion --
and the transformations are applied by the tool. Authoring it pre-factored would have hidden
exactly the work this component is graded on.

\ProductionCount\ productions over \TerminalCount\ terminals and \NonterminalCount\
nonterminals:

\begin{longtable}{@{}l l l@{}}
\toprule
\textbf{Nonterminal} & \textbf{Production} & \textbf{Rule id} \\
\midrule
\endhead
\bottomrule
\endfoot
\texttt{S} & $\rightarrow$ \texttt{Seq} & \scriptsize\texttt{demo.gram.s1} \\
\addlinespace
\texttt{Seq} & $\rightarrow$ \texttt{Seq CONJ Unit} & \scriptsize\texttt{demo.gram.seq1} \\
\texttt{} & $\mid$ \texttt{Seq SEP Unit} & \scriptsize\texttt{demo.gram.seq2} \\
\texttt{} & $\mid$ \texttt{Seq SEP CONJ Unit} & \scriptsize\texttt{demo.gram.seq3} \\
\texttt{} & $\mid$ \texttt{Unit} & \scriptsize\texttt{demo.gram.seq4} \\
\addlinespace
\texttt{Unit} & $\rightarrow$ \texttt{INTJ Clause} & \scriptsize\texttt{demo.gram.unit1} \\
\texttt{} & $\mid$ \texttt{INTJ} & \scriptsize\texttt{demo.gram.unit2} \\
\texttt{} & $\mid$ \texttt{FORMULA} & \scriptsize\texttt{demo.gram.unit3} \\
\texttt{} & $\mid$ \texttt{Clause} & \scriptsize\texttt{demo.gram.unit4} \\
\addlinespace
\texttt{Clause} & $\rightarrow$ \texttt{PP ClauseOpt} & \scriptsize\texttt{demo.gram.clause1} \\
\texttt{} & $\mid$ \texttt{NP ClauseTail} & \scriptsize\texttt{demo.gram.clause2} \\
\texttt{} & $\mid$ \texttt{Predicate} & \scriptsize\texttt{demo.gram.clause3} \\
\addlinespace
\texttt{ClauseOpt} & $\rightarrow$ \texttt{Clause} & \scriptsize\texttt{demo.gram.copt1} \\
\texttt{} & $\mid$ $\varepsilon$ & \scriptsize\texttt{demo.gram.copt2} \\
\addlinespace
\texttt{ClauseTail} & $\rightarrow$ \texttt{Predicate} & \scriptsize\texttt{demo.gram.ctail1} \\
\texttt{} & $\mid$ \texttt{NP Predicate} & \scriptsize\texttt{demo.gram.ctail2} \\
\texttt{} & $\mid$ $\varepsilon$ & \scriptsize\texttt{demo.gram.ctail3} \\
\addlinespace
\texttt{Predicate} & $\rightarrow$ \texttt{ADV Predicate} & \scriptsize\texttt{demo.gram.pred1} \\
\texttt{} & $\mid$ \texttt{NegOpt VerbGroup Comps} & \scriptsize\texttt{demo.gram.pred2} \\
\addlinespace
\texttt{NegOpt} & $\rightarrow$ \texttt{NEG} & \scriptsize\texttt{demo.gram.neg1} \\
\texttt{} & $\mid$ $\varepsilon$ & \scriptsize\texttt{demo.gram.neg2} \\
\addlinespace
\texttt{VerbGroup} & $\rightarrow$ \texttt{VERB VerbTail} & \scriptsize\texttt{demo.gram.vg1} \\
\addlinespace
\texttt{VerbTail} & $\rightarrow$ \texttt{VERB VerbTail} & \scriptsize\texttt{demo.gram.vt1} \\
\texttt{} & $\mid$ \texttt{ADV VerbTail} & \scriptsize\texttt{demo.gram.vt2} \\
\texttt{} & $\mid$ \texttt{NEG VerbTail} & \scriptsize\texttt{demo.gram.vt3} \\
\texttt{} & $\mid$ $\varepsilon$ & \scriptsize\texttt{demo.gram.vt4} \\
\addlinespace
\texttt{Comps} & $\rightarrow$ \texttt{NP Comps} & \scriptsize\texttt{demo.gram.comps1} \\
\texttt{} & $\mid$ \texttt{PREP CompObj Comps} & \scriptsize\texttt{demo.gram.comps2} \\
\texttt{} & $\mid$ \texttt{ADJ Comps} & \scriptsize\texttt{demo.gram.comps3} \\
\texttt{} & $\mid$ $\varepsilon$ & \scriptsize\texttt{demo.gram.comps4} \\
\addlinespace
\texttt{CompObj} & $\rightarrow$ \texttt{NP} & \scriptsize\texttt{demo.gram.cobj1} \\
\texttt{} & $\mid$ \texttt{VerbGroup} & \scriptsize\texttt{demo.gram.cobj2} \\
\addlinespace
\texttt{NP} & $\rightarrow$ \texttt{DET NOUN} & \scriptsize\texttt{demo.gram.np1} \\
\texttt{} & $\mid$ \texttt{DET ADJ NOUN} & \scriptsize\texttt{demo.gram.np2} \\
\texttt{} & $\mid$ \texttt{PRON} & \scriptsize\texttt{demo.gram.np3} \\
\texttt{} & $\mid$ \texttt{NOUN} & \scriptsize\texttt{demo.gram.np4} \\
\addlinespace
\texttt{PP} & $\rightarrow$ \texttt{PREP NP} & \scriptsize\texttt{demo.gram.pp1} \\
\addlinespace
\end{longtable}


The design decisions worth naming:

\begin{description}[leftmargin=1.4em, style=nextline]
  \item[\nonterm{Seq} is left recursive] \nonterm{Seq} $\rightarrow$ \nonterm{Seq}
    \term{CONJ} \nonterm{Unit} is the natural way to say ``a statement is a list of units
    joined by conjunctions or separators''. Three of its four alternatives recurse on the
    left, which is what the elimination step below removes.
  \item[\nonterm{Clause} has three entries] one for a fronted \nonterm{PP}, one starting
    from an \nonterm{NP}, and one for a bare \nonterm{Predicate} with no subject at all.
  \item[\nonterm{ClauseTail} allows two noun phrases] \nonterm{NP} \nonterm{Predicate}
    covers left-dislocated topics -- \fw{le taximan, lui il sait} -- which are extremely
    common and would otherwise be rejected.
  \item[\nonterm{Predicate} is recursive through \nonterm{ADV}] so adverbs stack freely
    before the verb group, which is where they actually occur.
  \item[\nonterm{VerbTail} is the serial-verb device] it lets verbs, adverbs and negations
    chain after the head verb without bounding the number of them.
\end{description}

\subsection{Transformation ledger}

The analyzer records every transformation it performs, with the identifiers of the rules
produced. \LedgerSteps\ steps were required: one left-recursion elimination and three left
factorings.

\begin{longtable}{@{}p{0.20\linewidth}p{0.46\linewidth}p{0.26\linewidth}@{}}
\toprule
\textbf{Transformation} & \textbf{Description} & \textbf{Rules produced} \\
\midrule
\endhead
\bottomrule
\endfoot
left recursion elimination & Eliminated direct left recursion on 'Seq' using fresh nonterminal "Seq'" & \scriptsize\texttt{demo.gram.seq4.lr, demo.gram.seq1.tail, demo.gram.seq2.tail, demo.gram.seq3.tail, Seq'.eps} \\
\addlinespace
left factoring & Factored common prefix ('SEP',) on "Seq'" using fresh nonterminal "Seq''" & \scriptsize\texttt{Seq'.factor0, demo.gram.seq2.tail.tail0, demo.gram.seq3.tail.tail0} \\
\addlinespace
left factoring & Factored common prefix ('INTJ',) on 'Unit' using fresh nonterminal "Unit'" & \scriptsize\texttt{Unit.factor1, demo.gram.unit1.tail1, demo.gram.unit2.tail1} \\
\addlinespace
left factoring & Factored common prefix ('DET',) on 'NP' using fresh nonterminal "NP'" & \scriptsize\texttt{NP.factor2, demo.gram.np1.tail2, demo.gram.np2.tail2} \\
\addlinespace
\end{longtable}


\subsubsection*{Step 1 --- left-recursion elimination on \nonterm{Seq}}

Applying the standard construction, the immediately left-recursive
\nonterm{Seq} $\rightarrow$ \nonterm{Seq}$\,\alpha_i \mid \beta$ becomes:

\begin{center}
\begin{tabular}{@{}l@{\quad}l@{}}
\nonterm{Seq}  & $\rightarrow$ \nonterm{Unit} \nonterm{Seq$'$} \\
\nonterm{Seq$'$} & $\rightarrow$ \term{CONJ} \nonterm{Unit} \nonterm{Seq$'$} \\
               & $\mid$ \term{SEP} \nonterm{Unit} \nonterm{Seq$'$} \\
               & $\mid$ \term{SEP} \term{CONJ} \nonterm{Unit} \nonterm{Seq$'$} \\
               & $\mid$ $\varepsilon$ \\
\end{tabular}
\end{center}

\subsubsection*{Steps 2--4 --- left factoring}

Elimination immediately \emph{creates} a new conflict: two \nonterm{Seq$'$} alternatives now
begin with \term{SEP}, so a predictive parser cannot choose between them on one token of
lookahead. That is factored into \nonterm{Seq$''$}. Two pre-existing conflicts are factored
the same way:

\begin{center}
\begin{tabular}{@{}lll@{}}
\toprule
\textbf{Nonterminal} & \textbf{Common prefix} & \textbf{Fresh nonterminal} \\
\midrule
\nonterm{Seq$'$} & \term{SEP} & \nonterm{Seq$''$} \\
\nonterm{Unit}   & \term{INTJ} & \nonterm{Unit$'$} \\
\nonterm{NP}     & \term{DET}  & \nonterm{NP$'$} \\
\bottomrule
\end{tabular}
\end{center}

\nonterm{NP} is the clearest illustration: \nonterm{NP} $\rightarrow$ \term{DET} \term{NOUN}
and \nonterm{NP} $\rightarrow$ \term{DET} \term{ADJ} \term{NOUN} share \term{DET}, so the
parser must delay the decision until after the determiner is consumed:

\begin{center}
\begin{tabular}{@{}l@{\quad}l@{}}
\nonterm{NP}   & $\rightarrow$ \term{DET} \nonterm{NP$'$} $\mid$ \term{PRON} $\mid$ \term{NOUN} \\
\nonterm{NP$'$} & $\rightarrow$ \term{NOUN} $\mid$ \term{ADJ} \term{NOUN} \\
\end{tabular}
\end{center}

\begin{keypoint}[Order matters]
Left factoring is applied \emph{after} left-recursion elimination, and repeatedly until no
common prefixes remain. Running them in the other order would leave the \term{SEP} conflict
that elimination introduces, and the grammar would not be LL(1).
\end{keypoint}

\subsection{Transformed grammar}

The executable grammar, \ExecProductionCount\ productions after transformation:

\input{generated/grammar-transformed}

\subsection{FIRST and FOLLOW sets}

FIRST and FOLLOW are computed by fixed-point iteration in
\texttt{core/ll1.py}: repeatedly apply the standard closure rules until no set changes.
Nullable nonterminals are those whose FIRST set derivation reaches $\varepsilon$; in this
grammar they are \nonterm{ClauseOpt}, \nonterm{ClauseTail}, \nonterm{Comps},
\nonterm{NegOpt}, \nonterm{Seq$'$}, \nonterm{Unit$'$} and \nonterm{VerbTail}.

\begin{longtable}{@{}l p{0.38\linewidth} p{0.38\linewidth}@{}}
\toprule
\textbf{Symbol} & \textbf{FIRST} & \textbf{FOLLOW} \\
\midrule
\endhead
\bottomrule
\endfoot
\texttt{S} & \small ADV, DET, FORMULA, INTJ, NEG, NOUN, PREP, PRON, VERB & \small \textbackslash{}\$ \\
\texttt{Seq} & \small ADV, DET, FORMULA, INTJ, NEG, NOUN, PREP, PRON, VERB & \small \textbackslash{}\$ \\
\texttt{Unit} & \small ADV, DET, FORMULA, INTJ, NEG, NOUN, PREP, PRON, VERB & \small \textbackslash{}\$, CONJ, SEP \\
\texttt{Clause} & \small ADV, DET, NEG, NOUN, PREP, PRON, VERB & \small \textbackslash{}\$, CONJ, SEP \\
\texttt{ClauseOpt} & \small ADV, DET, NEG, NOUN, PREP, PRON, VERB & \small \textbackslash{}\$, CONJ, SEP \\
\texttt{ClauseTail} & \small ADV, DET, NEG, NOUN, PRON, VERB & \small \textbackslash{}\$, CONJ, SEP \\
\texttt{Predicate} & \small ADV, NEG, VERB & \small \textbackslash{}\$, CONJ, SEP \\
\texttt{NegOpt} & \small NEG & \small VERB \\
\texttt{VerbGroup} & \small VERB & \small \textbackslash{}\$, ADJ, CONJ, DET, NOUN, PREP, PRON, SEP \\
\texttt{VerbTail} & \small ADV, NEG, VERB & \small \textbackslash{}\$, ADJ, CONJ, DET, NOUN, PREP, PRON, SEP \\
\texttt{Comps} & \small ADJ, DET, NOUN, PREP, PRON & \small \textbackslash{}\$, CONJ, SEP \\
\texttt{CompObj} & \small DET, NOUN, PRON, VERB & \small \textbackslash{}\$, ADJ, CONJ, DET, NOUN, PREP, PRON, SEP \\
\texttt{NP} & \small DET, NOUN, PRON & \small \textbackslash{}\$, ADJ, ADV, CONJ, DET, NEG, NOUN, PREP, PRON, SEP, VERB \\
\texttt{PP} & \small PREP & \small \textbackslash{}\$, ADV, CONJ, DET, NEG, NOUN, PREP, PRON, SEP, VERB \\
\texttt{Seq'} & \small CONJ, SEP & \small \textbackslash{}\$ \\
\texttt{Seq''} & \small ADV, CONJ, DET, FORMULA, INTJ, NEG, NOUN, PREP, PRON, VERB & \small \textbackslash{}\$ \\
\texttt{Unit'} & \small ADV, DET, NEG, NOUN, PREP, PRON, VERB & \small \textbackslash{}\$, CONJ, SEP \\
\texttt{NP'} & \small ADJ, NOUN & \small \textbackslash{}\$, ADJ, ADV, CONJ, DET, NEG, NOUN, PREP, PRON, SEP, VERB \\
\end{longtable}


\subsection{LL(1) parsing table}

The table is built by the standard rule: for each production
$A \rightarrow \alpha$, add it at $[A, t]$ for every $t \in \mathrm{FIRST}(\alpha)$; and if
$\alpha$ is nullable, also for every $t \in \mathrm{FOLLOW}(A)$. A second production landing
in an occupied cell is a conflict.

\begin{keypoint}[Result: conflict-free]
The constructed table contains \textbf{no duplicated cells}. The grammar is LL(1) and the
parser needs exactly one token of lookahead, with no backtracking and no ambiguity
resolution heuristics. \texttt{tests/test\_ll1.py} asserts this property, so a future
grammar edit that introduces a conflict fails the test suite rather than degrading
silently.
\end{keypoint}

The full table is on the next page; rows are nonterminals, columns are terminals plus the
end marker \term{\$}, and cells hold the production identifier to apply.

\begin{landscape}
\begin{longtable}{@{}l*{13}{c}@{}}
\toprule
\textbf{NT} & \rotatebox{90}{\texttt{DET}} & \rotatebox{90}{\texttt{NOUN}} & \rotatebox{90}{\texttt{PRON}} & \rotatebox{90}{\texttt{VERB}} & \rotatebox{90}{\texttt{PREP}} & \rotatebox{90}{\texttt{ADJ}} & \rotatebox{90}{\texttt{ADV}} & \rotatebox{90}{\texttt{NEG}} & \rotatebox{90}{\texttt{CONJ}} & \rotatebox{90}{\texttt{SEP}} & \rotatebox{90}{\texttt{INTJ}} & \rotatebox{90}{\texttt{FORMULA}} & \rotatebox{90}{\texttt{\$}} \\
\midrule
\endhead
\bottomrule
\endfoot
\texttt{S} & \tiny\texttt{s1} & \tiny\texttt{s1} & \tiny\texttt{s1} & \tiny\texttt{s1} & \tiny\texttt{s1} &  & \tiny\texttt{s1} & \tiny\texttt{s1} &  &  & \tiny\texttt{s1} & \tiny\texttt{s1} &  \\
\texttt{Seq} & \tiny\texttt{lr} & \tiny\texttt{lr} & \tiny\texttt{lr} & \tiny\texttt{lr} & \tiny\texttt{lr} &  & \tiny\texttt{lr} & \tiny\texttt{lr} &  &  & \tiny\texttt{lr} & \tiny\texttt{lr} &  \\
\texttt{Unit} & \tiny\texttt{unit4} & \tiny\texttt{unit4} & \tiny\texttt{unit4} & \tiny\texttt{unit4} & \tiny\texttt{unit4} &  & \tiny\texttt{unit4} & \tiny\texttt{unit4} &  &  & \tiny\texttt{factor1} & \tiny\texttt{unit3} &  \\
\texttt{Clause} & \tiny\texttt{clause2} & \tiny\texttt{clause2} & \tiny\texttt{clause2} & \tiny\texttt{clause3} & \tiny\texttt{clause1} &  & \tiny\texttt{clause3} & \tiny\texttt{clause3} &  &  &  &  &  \\
\texttt{ClauseOpt} & \tiny\texttt{copt1} & \tiny\texttt{copt1} & \tiny\texttt{copt1} & \tiny\texttt{copt1} & \tiny\texttt{copt1} &  & \tiny\texttt{copt1} & \tiny\texttt{copt1} & \tiny\texttt{copt2} & \tiny\texttt{copt2} &  &  & \tiny\texttt{copt2} \\
\texttt{ClauseTail} & \tiny\texttt{ctail2} & \tiny\texttt{ctail2} & \tiny\texttt{ctail2} & \tiny\texttt{ctail1} &  &  & \tiny\texttt{ctail1} & \tiny\texttt{ctail1} & \tiny\texttt{ctail3} & \tiny\texttt{ctail3} &  &  & \tiny\texttt{ctail3} \\
\texttt{Predicate} &  &  &  & \tiny\texttt{pred2} &  &  & \tiny\texttt{pred1} & \tiny\texttt{pred2} &  &  &  &  &  \\
\texttt{NegOpt} &  &  &  & \tiny\texttt{neg2} &  &  &  & \tiny\texttt{neg1} &  &  &  &  &  \\
\texttt{VerbGroup} &  &  &  & \tiny\texttt{vg1} &  &  &  &  &  &  &  &  &  \\
\texttt{VerbTail} & \tiny\texttt{vt4} & \tiny\texttt{vt4} & \tiny\texttt{vt4} & \tiny\texttt{vt1} & \tiny\texttt{vt4} & \tiny\texttt{vt4} & \tiny\texttt{vt2} & \tiny\texttt{vt3} & \tiny\texttt{vt4} & \tiny\texttt{vt4} &  &  & \tiny\texttt{vt4} \\
\texttt{Comps} & \tiny\texttt{comps1} & \tiny\texttt{comps1} & \tiny\texttt{comps1} &  & \tiny\texttt{comps2} & \tiny\texttt{comps3} &  &  & \tiny\texttt{comps4} & \tiny\texttt{comps4} &  &  & \tiny\texttt{comps4} \\
\texttt{CompObj} & \tiny\texttt{cobj1} & \tiny\texttt{cobj1} & \tiny\texttt{cobj1} & \tiny\texttt{cobj2} &  &  &  &  &  &  &  &  &  \\
\texttt{NP} & \tiny\texttt{factor2} & \tiny\texttt{np4} & \tiny\texttt{np3} &  &  &  &  &  &  &  &  &  &  \\
\texttt{PP} &  &  &  &  & \tiny\texttt{pp1} &  &  &  &  &  &  &  &  \\
\texttt{Seq'} &  &  &  &  &  &  &  &  & \tiny\texttt{tail} & \tiny\texttt{factor0} &  &  & \tiny\texttt{eps} \\
\texttt{Seq''} & \tiny\texttt{tail0} & \tiny\texttt{tail0} & \tiny\texttt{tail0} & \tiny\texttt{tail0} & \tiny\texttt{tail0} &  & \tiny\texttt{tail0} & \tiny\texttt{tail0} & \tiny\texttt{tail0} &  & \tiny\texttt{tail0} & \tiny\texttt{tail0} &  \\
\texttt{Unit'} & \tiny\texttt{tail1} & \tiny\texttt{tail1} & \tiny\texttt{tail1} & \tiny\texttt{tail1} & \tiny\texttt{tail1} &  & \tiny\texttt{tail1} & \tiny\texttt{tail1} & \tiny\texttt{tail1} & \tiny\texttt{tail1} &  &  & \tiny\texttt{tail1} \\
\texttt{NP'} &  & \tiny\texttt{tail2} &  &  &  & \tiny\texttt{tail2} &  &  &  &  &  &  &  \\
\end{longtable}

\end{landscape}
